% GENERATED FILE. Edit canonical lessons, then run npm run generate:books.
% canonical-source-hash: fnv1a64:a3f2b2ea6c5c2fb6
% canonical-lessons: MW-C287-sikhavno, MW-C287-bhanavno, MW-C287-guru, MW-C287-chelo, MW-C287-akharo

\chapter{Teachers and Pupils}
\label{ch:mw-teachers-and-pupils}
\hlchaptermodality{287}

\begin{chapteropening}
\textbf{By the end of this chapter:} \emph{I can say who teaches me and what I can learn.}

Complete the last lesson of chapter 287: I can say who teaches me and what I can learn.
\end{chapteropening}

\section[\texorpdfstring{\mw{सिखावणो} (sikhāvṇo)}{सिखावणो (sikhāvṇo)}]{\mw{सिखावणो} (sikhāvṇo) — to teach a skill}

\label{lesson:MW-C287-sikhavno}



\begin{quote}
Before the new one: say the Marwadi for a class at school, then the Marwadi for to mark schoolwork.
\end{quote}

\subsection*{You'll want to know: \mw{सिखावणो}}
\textbf{\mw{सिखावणो}} — \emph{sikhāvṇo} — \textquotedblleft{}to teach a skill\textquotedblright{}.

\begin{quote}
\mw{दादी} \mw{सीवणो} \mw{सिखावै} \mw{है।}
\end{quote}

\begin{quote}
\emph{dādī sīvṇo sikhāvai hai.}
\end{quote}

\begin{quote}
Grandmother teaches sewing.
\end{quote}

Past (I …): \mw{म्हैं} \mw{गीत} \mw{सिखायो}, \emph{mhaĩ gīt sikhāyo}.

\begin{cousinweb}
\textbf{\mw{सिखावणो}} (\emph{sikhāvṇo}) is built on \textbf{\mw{सीखणो}}, to learn, as its causative: \textquotedblleft{}to make learn\textquotedblright{}. CDIAL gives no Marwari form in this entry, but lists Hindi \emph{sikhānā}, \textquotedblleft{}teaches\textquotedblright{}, under Sanskrit \emph{śikṣate} (CDIAL 12430).
\end{cousinweb}

\begin{grammarlens}[title={What you've built}]
One more word for this chapter.
\end{grammarlens}

\subsection*{Your turn}
\begin{itemize}
\raggedright
  \item \emph{Say it:} \emph{sikhāvṇo}
  \item \emph{Say it:} \emph{sikhāvṇo}, once more
  \item \emph{Say it:} \emph{dādī sīvṇo sikhāvai hai.} — Grandmother teaches sewing.
  \item \emph{Recall:} say the Marwadi for a class at school, then the Marwadi for to mark schoolwork, then say \emph{sikhāvṇo} again
\end{itemize}

\begin{tcolorbox}[breakable,colback=teal!4,colframe=teal!35!black,title={Before you move on}]
What does \mw{सिखावणो} mean? (\textquotedblleft{}To teach a skill\textquotedblright{}.) Say it once more.
\end{tcolorbox}
\section[\texorpdfstring{\mw{भणावणो} (bhaṇāvṇo)}{भणावणो (bhaṇāvṇo)}]{\mw{भणावणो} (bhaṇāvṇo) — to teach lessons}

\label{lesson:MW-C287-bhanavno}



\begin{quote}
Before the new one: say the Marwadi for to mark schoolwork, then the Marwadi for to teach a skill.
\end{quote}

\subsection*{You'll want to know: \mw{भणावणो}}
\textbf{\mw{भणावणो}} — \emph{bhaṇāvṇo} — \textquotedblleft{}to teach lessons\textquotedblright{}.

\begin{quote}
\mw{मास्टर} \mw{भासा} \mw{भणावै} \mw{है।}
\end{quote}

\begin{quote}
\emph{māsṭar bhāsā bhaṇāvai hai.}
\end{quote}

\begin{quote}
The teacher teaches language.
\end{quote}

Past (I …): \mw{म्हैं} \mw{पाठ} \mw{भणायो}, \emph{mhaĩ pāṭh bhaṇāyo}.

\begin{cousinweb}
\textbf{\mw{भणावणो}} (\emph{bhaṇāvṇo}) is built on \textbf{\mw{भणणो}}, to study, as its causative: \textquotedblleft{}to make study\textquotedblright{}. CDIAL lists Old Marwari \emph{bhaṇaï}, \textquotedblleft{}recites\textquotedblright{} (CDIAL 9383).
\end{cousinweb}

\begin{grammarlens}[title={What you've built}]
One more word for this chapter.
\end{grammarlens}

\subsection*{Your turn}
\begin{itemize}
\raggedright
  \item \emph{Say it:} \emph{bhaṇāvṇo}
  \item \emph{Say it:} \emph{bhaṇāvṇo}, once more
  \item \emph{Say it:} \emph{māsṭar bhāsā bhaṇāvai hai.} — The teacher teaches language.
  \item \emph{Recall:} say the Marwadi for to mark schoolwork, then the Marwadi for to teach a skill, then say \emph{bhaṇāvṇo} again
\end{itemize}

\begin{tcolorbox}[breakable,colback=teal!4,colframe=teal!35!black,title={Before you move on}]
What does \mw{भणावणो} mean? (\textquotedblleft{}To teach lessons\textquotedblright{}.) Say it once more.
\end{tcolorbox}
\section[\texorpdfstring{\mw{गुरु} (guru)}{गुरु (guru)}]{\mw{गुरु} (guru) — a guru; a master teacher}

\label{lesson:MW-C287-guru}



\begin{quote}
Before the new one: say the Marwadi for to teach a skill, then the Marwadi for to teach lessons.
\end{quote}

\subsection*{You'll want to know: \mw{गुरु}}
\textbf{\mw{गुरु}} — \emph{guru} — \textquotedblleft{}a guru; a master teacher\textquotedblright{}.

\begin{quote}
\mw{गुरु} \mw{संगीत} \mw{सिखावै} \mw{है।}
\end{quote}

\begin{quote}
\emph{guru saṅgīt sikhāvai hai.}
\end{quote}

\begin{quote}
The guru teaches music.
\end{quote}

A masculine noun.

\begin{cousinweb}
\textbf{\mw{गुरु}} (\emph{guru}) is a tatsama, taken from Sanskrit \emph{guru}, \textquotedblleft{}heavy\textquotedblright{}, and so \textquotedblleft{}venerable\textquotedblright{}; Monier-Williams also gives it as a teacher (Monier-Williams).
\end{cousinweb}

\begin{grammarlens}[title={What you've built}]
One more word for this chapter.
\end{grammarlens}

\subsection*{Your turn}
\begin{itemize}
\raggedright
  \item \emph{Say it:} \emph{guru}
  \item \emph{Say it:} \emph{guru}, once more
  \item \emph{Say it:} \emph{guru saṅgīt sikhāvai hai.} — The guru teaches music.
  \item \emph{Recall:} say the Marwadi for to teach a skill, then the Marwadi for to teach lessons, then say \emph{guru} again
\end{itemize}

\begin{tcolorbox}[breakable,colback=teal!4,colframe=teal!35!black,title={Before you move on}]
What does \mw{गुरु} mean? (\textquotedblleft{}A guru; a master teacher\textquotedblright{}.) Say it once more.
\end{tcolorbox}
\section[\texorpdfstring{\mw{चेलो} (chelo)}{चेलो (chelo)}]{\mw{चेलो} (chelo) — a pupil}

\label{lesson:MW-C287-chelo}



\begin{quote}
Before the new one: say the Marwadi for to teach lessons, then the Marwadi for a guru.
\end{quote}

\subsection*{You'll want to know: \mw{चेलो}}
\textbf{\mw{चेलो}} — \emph{chelo} — \textquotedblleft{}a pupil\textquotedblright{}.

\begin{quote}
\mw{यो} \mw{म्हारो} \mw{चेलो} \mw{है।}
\end{quote}

\begin{quote}
\emph{yo mhāro chelo hai.}
\end{quote}

\begin{quote}
This is my pupil.
\end{quote}

A masculine noun.

\begin{cousinweb}
\textbf{\mw{चेलो}} (\emph{chelo}) is inherited from a reconstructed \emph{*cēlla}, \textquotedblleft{}boy\textquotedblright{}; CDIAL gives no Marwari form in this entry, but lists Hindi \emph{celā}, \textquotedblleft{}pupil, apprentice\textquotedblright{}, Gujarati \emph{celo}, \textquotedblleft{}pupil\textquotedblright{}, and Sindhi \emph{celo}, \textquotedblleft{}disciple\textquotedblright{} (CDIAL 4911).
\end{cousinweb}

\begin{grammarlens}[title={What you've built}]
One more word for this chapter.
\end{grammarlens}

\subsection*{Your turn}
\begin{itemize}
\raggedright
  \item \emph{Say it:} \emph{chelo}
  \item \emph{Say it:} \emph{chelo}, once more
  \item \emph{Say it:} \emph{yo mhāro chelo hai.} — This is my pupil.
  \item \emph{Recall:} say the Marwadi for to teach lessons, then the Marwadi for a guru, then say \emph{chelo} again
\end{itemize}

\begin{tcolorbox}[breakable,colback=teal!4,colframe=teal!35!black,title={Before you move on}]
What does \mw{चेलो} mean? (\textquotedblleft{}A pupil\textquotedblright{}.) Say it once more.
\end{tcolorbox}
\section[\texorpdfstring{\mw{अखाड़ो} (akhāṛo)}{अखाड़ो (akhāṛo)}]{\mw{अखाड़ो} (akhāṛo) — a wrestling ground; a training ground}

\label{lesson:MW-C287-akharo}



\begin{quote}
Before the new one: say the Marwadi for a guru, then the Marwadi for a pupil.
\end{quote}

\subsection*{You'll want to know: \mw{अखाड़ो}}
\textbf{\mw{अखाड़ो}} — \emph{akhāṛo} — \textquotedblleft{}a wrestling ground; a training ground\textquotedblright{}.

\begin{quote}
\mw{गांव} \mw{मांय} \mw{एक} \mw{अखाड़ो} \mw{है।}
\end{quote}

\begin{quote}
\emph{gā̃v mā̃y ek akhāṛo hai.}
\end{quote}

\begin{quote}
There's a wrestling ground in the village.
\end{quote}

A masculine noun.

\begin{cousinweb}
\textbf{\mw{अखाड़ो}} (\emph{akhāṛo}) is inherited from Sanskrit \emph{akṣavāṭa}, \textquotedblleft{}a wrestling ground\textquotedblright{}; CDIAL gives no Marwari form in this entry, but lists Hindi \emph{akhāṛā}, \textquotedblleft{}wrestling ground\textquotedblright{}, and Gujarati \emph{akhāṛo}, \textquotedblleft{}wrestling place\textquotedblright{} (CDIAL 39).
\end{cousinweb}

\begin{grammarlens}[title={What you've built}]
One more word for this chapter.
\end{grammarlens}

\subsection*{Your turn}
\begin{itemize}
\raggedright
  \item \emph{Say it:} \emph{akhāṛo}
  \item \emph{Say it:} \emph{akhāṛo}, once more
  \item \emph{Say it:} \emph{gā̃v mā̃y ek akhāṛo hai.} — There's a wrestling ground in the village.
  \item \emph{Recall:} say the Marwadi for a guru, then the Marwadi for a pupil, then say \emph{akhāṛo} again
\end{itemize}

\begin{tcolorbox}[breakable,colback=teal!4,colframe=teal!35!black,title={Before you move on}]
What does \mw{अखाड़ो} mean? (\textquotedblleft{}A wrestling ground; a training ground\textquotedblright{}.) Say it once more.
\end{tcolorbox}
